\section{Matrix elements}

The nucleon spin-averaged matrix elements for operators composed of two-gluon-fields (with all
four indices non-contracted) are specified by
\begin{align}
{ M}_{\mu \alpha; \lambda \beta } (z,p) \equiv \langle p | \,
G_{\mu \alpha} (z) \, [z,0]\, G_{\lambda \beta} (0) | p \rangle \ ,
\end{align}
where $[z,0]$ is the standard straight-line gauge link
in the gluon (adjoint)
representation
\begin{align}
[x,y]~\equiv~{\rm Pexp}\Big\{ig\!\int_0^1\! dt~(x-y)^\mu \tilde{A}_\mu(tx+(1-t)y)\Big\}
\ .
\label{straightE}
\end{align}
The tensor structures for a decomposition over invariant amplitudes
may be built from two available 4-vectors $p_\alpha$, $z_\alpha$ and the
metric tensor $g_{\alpha \beta}$. Incorporating the antisymmetry of $G_{\rho \sigma}$ with respect to its indices, we have
\begin{align}
& { M}_{\mu \alpha; \lambda \beta} (z,p)= \nonumber \\ &
\left( g_{\mu\lambda} p_\alpha p_\beta - g_{\mu\beta} p_\alpha p_\lambda - g_{\alpha\lambda} p_\mu p_\beta + g_{\alpha\beta} p_\mu p_\lambda \right) \mathcal{M}_{pp} \nonumber \\
+& \left( g_{\mu\lambda} z_\alpha z_\beta - g_{\mu\beta} z_\alpha z_\lambda - g_{\alpha\lambda} z_\mu z_\beta + g_{\alpha\beta} z_\mu z_\lambda \right) \mathcal{M}_{zz} \nonumber \\
+ & \left( g_{\mu\lambda} z_\alpha p_\beta - g_{\mu\beta} z_\alpha p_\lambda - g_{\alpha\lambda} z_\mu p_\beta + g_{\alpha\beta} z_\mu p_\lambda \right) \mathcal{M}_{zp} \nonumber \\
+ & \left( g_{\mu\lambda} p_\alpha z_\beta - g_{\mu\beta} p_\alpha z_\lambda - g_{\alpha\lambda} p_\mu z_\beta + g_{\alpha\beta} p_\mu z_\lambda \right) \mathcal{M}_{pz} \nonumber \\
+& \left( p_\mu z_\alpha - p_\alpha z_\mu\right) \left( p_\lambda z_\beta - p_\beta z_\lambda\right) \mathcal{M}_{ppzz}
\nn + & \left(g_{\mu\lambda} g_{\alpha\beta} -g_{\mu\beta} g_{\alpha\lambda} \right)\mathcal{M}_{gg} \ ,
\label{Manb}
\end{align}
where the amplitudes ${\cal M}$ are functions of
the invariant interval $z^2$ and the Ioffe time \cite{Ioffe:1969kf}
$(pz)\equiv - \nu$ (the minus sign
is introduced for further convenience).

Since the matrix element should be symmetric with respect to interchange of the
fields (which amounts to $\{\mu \alpha\} \leftrightarrow \{\lambda \beta\}$ and $z \to -z$),
the functions $\mathcal{M}_{pp}$, $\mathcal{M}_{zz}$, $\mathcal{M}_{gg}$, $\mathcal{M}_{ppzz}$ and
$\mathcal{M}_{pz} - \mathcal{M}_{zp}$ are
even functions of $\nu$, while $\mathcal{M}_{pz} + \mathcal{M}_{zp}$ is
odd in $\nu$.

The usual light-cone gluon distribution is obtained from
$g^{\alpha \beta} { M}_{+ \alpha; \beta +} (z,p)$, with $z$ taken in the light-cone ``minus'' direction,
$z=z_-$. We have
\begin{align}
g^{\alpha \beta} { M}_{+ \alpha; \beta +} (z_-,p) =-2 p_+ ^2 \mathcal{M}_{pp}(\nu ,0) \ ,
\label{lcpp}
\end{align}
i.e., the PDF is determined by the ${\cal M}_{pp}$ structure,
\begin{align}
- \mathcal{M}_{pp} (\nu,0) =\frac12
\int_{-1}^1 \dd x \, e^{-i x \nu} x f_g (x) \, \ .
\label{glPDF}
\end{align}
Thus, we should choose the operators with the sets
$\{\mu \alpha; \lambda \beta \}$ that contain ${\cal M}_{pp}$
in their parametrization.

Note that it is the density of the momentum \mbox{$G(x) \equiv x f_g (x)$}
carried by the gluons rather than their number
density $f_g (x)$ that is a natural quantity in this definition of the gluon PDF.
In the local $z_-=0$ (or $\nu=0$) limit, the $x$-integral gives the fraction of the
hadron's plus momentum carried by the gluons.
In the absence of gluon-quark transitions, this fraction is conserved,
which puts a restriction on the $gg$-component of the Altarelli-Parisi \cite{Altarelli:1977zs} kernel.
Namely, it should have the plus-prescription
property when applied to $G(x)$.

Due to antisymmetry of $G_{\rho \sigma}$ with respect to its indices,
the values $\alpha=+$ and $\beta= +$ are excluded from the summation in Eq. (\ref{lcpp}).
Furthermore, since $g_{- -} =0$, the combination
$g^{\alpha \beta} { M}_{+ \alpha; \beta +} (z,p)$ includes only summation over
transverse indices $i,j =1,2$, i.e. reduces
to $g^{ij} { M}_{+ i; j +} (z,p) \equiv { M}_{+ i; + i} (z,p)$
(we switched here to Euclidean summation over $i$), for which we have
\begin{align}
{ M}_{+ i; +i}
= { M}_{0 i; 0i} + { M}_{3 i; 3i} + ({ M}_{0 i; 3i} +
{ M}_{3 i; 0i}) \ .
\label{ii}
\end{align}
In the local $z_3=0$ limit, these three combinations are proportional to ${\bf E}^2_\perp$, ${\bf B}^2_\perp$
and the third component $({\bf E \times B})_3$ of the Poynting vector,
respectively.

The decomposition of these combinations (with summation over $i$) in the basis of the ${\cal M}$ structures is
\begin{align}
{ M}_{0 i; i 0 } = & 2 p_0^2 \mathcal{M}_{pp} +2 \mathcal{M}_{gg} \ , \\
{ M}_{3 i; i 3 } = &2 p_3^2 \mathcal{M}_{pp}+2 z_3^2 \mathcal{M}_{zz} \nonumber \\ &
+2 z_3 p_3 \left( \mathcal{M}_{zp} + \mathcal{M}_{pz} \right) - 2\mathcal{M}_{gg} \ ,
\\
{ M}_{0 i; i 3 } = & 2 p_0 \left( p_3 \mathcal{M}_{pp} + z_3 \mathcal{M}_{pz} \right) \ ,
\\
{ M}_{3 i; i 0 } = & 2 p_0 \left( p_3 \mathcal{M}_{pp} + z_3 \mathcal{M}_{zp} \right) \ .
\end{align}

All of them contain the $\mathcal{M}_{pp}$ function defining the gluon
distribution, though with different kinematical factors.
Unfortunately, none of them is just $\mathcal{M}_{pp}$: they all contain
contaminating terms.
Moreover, the ${ M}_{3 i; i 3}$ matrix element (proposed
originally \cite{Ji:2013dva} for extractions of the gluon PDF on the lattice)
contains three contaminations, while the others have
just one addition.
In particular, the matrix element ${ M}_{0 i; i 0}$ has $\mathcal{M}_{gg}$ as a contaminating term.
It is easy to see that
\begin{align}
{ M}_{j i ; ij} \equiv \bra{p} G_{j i} (z) G_{ij}(0) \ket{p} &= -2 \mathcal{M}_{gg} \ ,
\end{align}
where the summation over both $i$ and $j$ is assumed. Hence, the combination
\begin{align}
{ M}_{0 i; i 0} + { M}_{j i ; i j } = & 2 p_0^2 \mathcal{M}_{pp} \
\label{00m}
\end{align}
may be used for extraction
of the twist-2 function $\mathcal{M}_{pp}$.

Combining together matrix elements of different types, one should
take into account
that,
off the light cone, these matrix elements have extra ultraviolet divergences
related to presence of the gauge link.
Due to the local nature of ultraviolet divergences,
each matrix element, for any set of its indices $\{\mu \alpha; \lambda \beta\}$,
is multiplicatively renormalizable with respect to these divergences \cite{Li:2018tpe}.
However, choosing different sets of $\{\mu \alpha; \nu \beta\}$, we get, in general, different
anomalous dimensions.

Thus, it is not evident {\it a priori} which linear combinations of these
matrix elements are multiplicatively renormalizable.
In Ref. \cite{Zhang:2018diq}, it was established that the combinations represented in Eq. (\ref{ii}), namely,
${ M}_{0 i; i 0}$, ${ M}_{3 i; i 3}$, \mbox{${ M}_{0 i; i3} +
{ M}_{3 i; i 0}$} (and also \mbox{${ M}_{0 i; i 3} -
{ M}_{3 i; i 0}$}), with
summation over transverse indices $i$,
are each multiplicatively renormalizable at the one-loop level.

Furthermore, the combination $G_{ij}G_{ij}$ (with summation over transverse $i,j$)
equals to $2G_{12} G_{12}$,
whose matrix elements are multiplicatively renormalizable.
As we will see, it has the same one-loop UV anomalous dimension
as ${M}_{0 i; i 0}$, hence the combination of Eq. (\ref{00m})
is multiplicatively renormalizable at the one-loop level.
A possible subject for further studies is to investigate
if this is true in higher orders.

The combination $g^{\alpha \beta} { M}_{3 \alpha; 3 \beta}$,
containing a covariant summation over $\alpha$ and $\beta$,
was also found to be multiplicatively renormalizable.
It is given by
\begin{align}
& g^{\alpha \beta} { M}_{3 \alpha; 3 \beta } =
\left(2p_3^2 - m^2 \right) \mathcal{M}_{pp} + 3z_3^2 \mathcal{M}_{zz}
\nonumber \\ & + 3 p_3 z_3 \left( \mathcal{M}_{zp}+ \mathcal{M}_{pz} \right) + p_0^2 z_3^2 \mathcal{M}_{ppzz}
-3\mathcal{M}_{gg} \ ,
\end{align}
and has the largest number (four) of contaminations.

The function $g^{\alpha \beta} { M}_{0 \alpha; 0 \beta}$, also
involving a covariant summation, was used in the first attempt
\cite{Fan:2018dxu}
of the lattice extraction of the gluon PDF. However, as noted in Ref.
\cite{Zhang:2018diq}, it is not multiplicatively renormalizable.

In any theory with a dimensionless coupling constant,
the matrix elements ${ M}(z,p)$
contain \mbox{$\sim \ln (-z^2)$} terms
corresponding to perturbative (or ``DGLAP'' for Dokshitzer-Gribov-Lipatov-Altarelli-Parisi
\cite{Gribov:1972ri,Altarelli:1977zs,Dokshitzer:1977sg})
evolution. One may wonder which combinations have
a diagonal DGLAP evolution at one loop.

To answer these questions, we have
calculated the modification of the original bilocal operator by
one-loop
gluon exchanges.
